\chapter{Wprowadzenie do tematyki}
\section{Czujniki klimatyczne}
Współczesny rozwój technologiczny w~połączeniu z~postępującymi
zamianami klimatu oraz rosnącą automatyzacją procesów 
przemysłowych i~urządzeń domowych sprawiły, że precyzyjne
obserwowanie i~dokumentowanie parametrów środowiskowych stało się
nieodłącznym elementem współczesnej inżynierii.

Głównym zadaniem czujników klimatycznych jest ciągła 
i~bezobsługowa akwizycja danych dotyczących parametrów atmosfery 
w~wybranym obszarze. Dane te są kluczowe w~wielu dziedzinach
życia gospodarczego i~społecznego. W~rolnictwie precyzyjnym
\translation{smart agriculture} pozwalają na optymalizację 
nawadniania i~ochrony upraw [TODO: citation needed], 
w~aglomeracjach miejskich służą do monitorowania smogu i~poziomu
zanieczyszczeń powietrza, natomiast w~systemach zarządzania 
budynkami \translation{Building Management Systems} stanowią 
podstawę sterowania procesami klimatyzacji, wentylacji 
i~ogrzewania.

[TODO: Może jeszcze jeden akapit? Praca Wołodko? Modularność?]

\section{Moduły pomiarowe}
Do podstawowych wielkości fizycznych rejestrowanych przez 
czujniki klimatyczne zalicza się przede wszystkim temperaturę, 
wilgotność względną oraz ciśnienie atmosferyczne. Do pomiaru tych 
wielkości coraz rzadziej wykorzystuje sie bezpośrednio połączone
z~układem czujniki analogowe, a~są one wypierane przez moduły 
integrujące elementy pomiarowe, dedykowane przetworniki
analogowo-cyfrowe (ADC, ang. \textit{analog-digital converter})
oraz fabryczne obwody kalibracyjne w~jednej obudowie. Pozwala to, 
między innymi, wyeliminować problem zakłóceń elektromagnetycznych 
indukowanych w~przewodach sygnałowych. Ponadto, projektanci modułów
implementują algorytmy linearyzacji charakterystyk sensorów
oraz filtry szumów analogowych po stronie samego modułu, pozwalając
na wygodniejszych odczyt danych.


\section{Zasilanie czujników}
[Panele słoneczne i magazyn energii]

\chapter{Programowanie systemów wbudowanych}
\section{Języki programowania}
[C, C++, Rust, MicroPython, co do jakiego mikrokontrolera]
\section{Obsługa magistrali}
[Opisz czym jest I2C z perspektywy kodu]